% !TEX root = ../main.tex
% !TEX program = xelatex
\title{基于多源预测与线性规划的微网协同调度}
\date{}
\maketitle
\begingroup
\setstretch{1.22}
\begin{abstract}
\noindent\hspace*{2em}随着我国碳达峰、碳中和目标的推进和分布式新能源的发展，光伏与储能在社区供能中的应用不断拓展。然而，光伏出力与用电需求存在时间错位，提前订电又面临供需和价格变化：订电不足可能产生高价紧急购电，订电过多则可能付费后无法充分利用。为降低微网购电费用，本文以10分钟为基本时段，建立统一的电量流与储能状态关系，研究不同信息条件下的购电及充放电策略。

\textbf{针对问题一}，将全天购电与储能调度表述为\textbf{电量流线性规划问题}。在日初与日末储电量相等的条件下，统筹外网购电、光伏利用和电池充放电，将有限储能用于替代更高价时段的外网购电。求得代表日最低购电费用为\textbf{35126.85元}，较峰谷策略降低\textbf{8.23\%}。结果表明，结合全天供需变化进行\textbf{储能统筹安排}，能够避免过早充满或放空电池，进一步降低购电费用。

\textbf{针对问题二}，将\textbf{近期加权准备}与\textbf{跨日线性规划}相结合。每天零点依据已有历史，分别利用负载的星期规律和光伏的近期变化预测供需，再根据近期缺口误差确定购电准备量，平衡\textbf{五倍电价补购}与\textbf{过量订电}的损失。实际剩余储电量连续传递至次日，使各日安排相互衔接。2025年2月至12月总费用为\textbf{1358.38万元}，较历史误差等权处理的准备方案节省\textbf{44.66万元}，说明\textbf{偏重近期误差}有助于准备量跟随供需变化。

\textbf{针对问题三}，通过\textbf{光伏融合预测、日内负载校正与滚动线性规划}更新购电安排。利用当天已发生的负载偏差修正后续用电判断，并融合官方光伏预报与历史参考，综合调整费用和当前储电量，重新优化尚未执行的订单。同一评价期内总费用为\textbf{1301.42万元}。以零点固定订单为对照，日内运行反馈与订单调整节省\textbf{50.63万元}；在相同调整权限下，引入后发官方光伏预报进一步节省\textbf{6.95万元}。模型将预报更新与调整费用共同纳入决策，在取得额外节费的同时，使\textbf{紧急购电量下降30.10\%}，体现了日内新预报的\textbf{经济与运行价值}。

\textbf{针对问题四}，在未来实际电价不可提前获知的条件下，为两条分支分别配置价格预测方法。零点分支以\textbf{提升树预测与线性规划}连接价格判断和购电安排，刻画历史价格与供需特征之间的非线性关系；日内分支借助\textbf{岭回归预测与滚动线性规划}，处理相关价格特征，并利用当天已发生的价格修正后续判断。预测电价用于决策，实际电价用于结算。两套完整策略同期总费用分别为\textbf{1424.86万元}和\textbf{1369.59万元}，\textbf{日内策略降低3.88\%}，体现了供需与价格判断及时更新、购电计划随之调整的综合价值。

逐时供需与费用复核表明，所得方案满足供需平衡及储能容量、功率约束。准备量\textbf{灵敏度分析}显示，准备不足会增加紧急购电，过度准备则增加普通订电支出，\textbf{适度准备更为经济}。本文将预测误差、储能调节与购电费用联系起来，为社区微网在不同信息条件下的电力调控提供了可解释的决策方法。
\end{abstract}
\noindent{\heiti\zihao{-4} 关键词：}{\songti\zihao{-4}\textbf{微网调度；储能优化；线性规划；滚动优化；电价预测}}
\par
\endgroup
